\section{Identification protocol and provenance}
\label{sec:identification-details}

\subsection{Exact-range control at $151.9$M}

\paragraph{What is pinned.} The control is designed so that every scalar
quantity a base/range method can express is identical between arms. Within each
seed pair the two arms share: model architecture ($151{,}898{,}880$
parameters); trainable initialisation; training length $256$; token order;
optimiser, LR schedule, global batch; within-pair micro batch and accumulation;
training budget
($499{,}974{,}144$ tokens, $7{,}629$ optimiser steps); the $32$ frozen
evaluation anchors; the sampled highest frequency; the sampled lowest
frequency; and the log-frequency span. The single free variable is the location
of the $K{-}2 = 30$ interior frequencies.

The optimiser is fused AdamW with $\beta_1{=}0.9$, $\beta_2{=}0.95$, weight
decay $0.01$, and gradient clipping at $1.0$. Peak/minimum learning rates are
$6\times10^{-4}/6\times10^{-5}$ with $762$ warmup steps and cosine decay.
The global batch is $256$ sequences of length $256$ for every arm. Seed $42$
uses micro batch $64$ with four accumulation steps; seeds $137$ and $256$ use
micro batch $128$ with two accumulation steps. Each paired comparison therefore
has identical execution geometry, optimiser steps, and counted-token budget;
the aggregate uses training seed as its statistical unit under this shared
scientific protocol.

\paragraph{Arms.} \emph{Baseline:} a paper-faithful reimplementation of the
uniform-in-log exponent grid specified in \S6.1 of \citet{oka2026fmrope}, at
training base $256$. The baseline implements that published rule directly.
\emph{Intervention:}
endpoint-normalised cosh interior spacing at $\tau{=}4$. Note that the
endpoint-normalised arm and deployed midpoint \evq{} are different support
embeddings of the \emph{same normalised shape}: both use midpoint Cosh
quantiles $q_k$, and the diagnostic removes only their affine endpoint degrees
of freedom. This normalisation is what makes the exact-range control possible.

Concretely, let $q_k{=}\Phi_\tau((k+\tfrac12)/K)$,
$s_k{=}(q_k-q_0)/(q_{K-1}-q_0)$, and
$R{=}((K-1)/K)\log B$.  The intervention uses
\begin{equation}
\omega_k^{\mathrm{anchor}}=\exp(-R s_k),
\qquad
\omega_k^{\mathrm{Geo}}=\exp\!\left(-R\frac{k}{K-1}\right).
\label{eq:exact-range-grid}
\end{equation}
Hence $k{=}0,K{-}1$ match exactly and only the $K{-}2$ interior entries differ.

\paragraph{Evaluation.} Two conditions. \texttt{fixed}: each checkpoint retains
its training range. \texttt{target-matched}: the range is set to the declared
evaluation length while preserving each arm's normalised interior spacing.
Metric is paired final-$128$-token teacher-forced NLL over the $32$ frozen
anchors; reported PPL is $\exp(\mathrm{mean\ NLL})$.

\paragraph{Three-seed contrasts.}
Negative values favour Cosh. The aggregation unit is the training seed;
evaluation anchors remain paired observations within a seed.

\begin{table}[ht]
\caption{Exact-range Cosh-minus-uniform tail-NLL contrasts. Mean$\pm$SD is
across three independent training seeds; no significance claim is attached to
$n{=}3$.}
\label{tab:exact-range-3seed}
\centering
\small
\resizebox{\linewidth}{!}{%
\begin{tabular}{@{}r rrr c rrr c@{}}
\toprule
& \multicolumn{4}{c}{Fixed training range}
& \multicolumn{4}{c}{Target-matched range} \\
Length & Seed 42 & Seed 137 & Seed 256 & Mean$\pm$SD
       & Seed 42 & Seed 137 & Seed 256 & Mean$\pm$SD \\
\midrule
$256$  & $+0.033$ & $+0.029$ & $+0.017$ & $+0.026\pm0.008$
       & $+0.033$ & $+0.029$ & $+0.017$ & $+0.026\pm0.008$ \\
$512$  & $-0.478$ & $-0.271$ & $-0.094$ & $-0.281\pm0.192$
       & $+0.061$ & $+0.053$ & $+0.067$ & $+0.060\pm0.007$ \\
$1{,}024$ & $-0.205$ & $-0.191$ & $-0.132$ & $-0.176\pm0.039$
       & $+0.182$ & $+0.156$ & $+0.344$ & $+0.227\pm0.102$ \\
$2{,}048$ & $-0.113$ & $-0.181$ & $-0.143$ & $-0.146\pm0.034$
       & $+0.279$ & $+0.399$ & $+0.701$ & $+0.460\pm0.218$ \\
\bottomrule
\end{tabular}}
\end{table}

Under \texttt{fixed}, Cosh is favoured by $3/3$ seeds at every OOD length;
under \texttt{target-matched}, uniform FMRoPE is favoured by $3/3$ seeds at
every OOD length. The fixed-range sign agrees across all seeds and OOD lengths;
the $1$K/$2$K magnitudes are the most stable across seeds.

\paragraph{Reading.} The \texttt{fixed} column identifies the optimisation
variable: interior allocation is not reducible to scalar base/range selection.
The \texttt{target-matched} columns measure the deployment interaction when
support is retargeted for both arms. Support and allocation are separately
controllable and jointly empirical coordinates (\S\ref{sec:exp-identify}).

\subsection{Exact-range factorial at $50.9$M}

\paragraph{Design.} The schedule arms, including the $1.25\times$ multiplier,
were specified before any result was inspected. Grid: $2$ bases $\times$
$2$ training lengths $\times$ $3$ head dimensions $=12$ structural
configurations; $5$ schedules (Geo, anchored cosh at $0.75\times$, $1.00\times$
[the formula point], $1.25\times$, and a deformation-matched exponential);
$3$ seeds. Main factorial $180/180$ completed, boundary-$\tau$ follow-up
$12/12$ completed, covering the complete pre-specified matrix.

\paragraph{Metric.} Weighted OOD NLL over $2L/4L/8L$,
\[
\bar N_{\mathrm{OOD}}
= \frac{\sum_{r\in\{2,4,8\}}\log_2(r+1)\,N_{rL}}
       {\sum_{r\in\{2,4,8\}}\log_2(r+1)},
\]
paired within structural configuration and seed, averaged within seed, then
equally weighted over the $12$ configurations. The statistical unit is the
configuration. Equivalent-PPL figures quoted anywhere are
$\exp(\mathrm{mean\ NLL})$, never an arithmetic mean of raw PPL.

\begin{table}[tb]
\centering
\small
\setlength{\tabcolsep}{2.5pt}
\renewcommand{\arraystretch}{0.92}
\caption{\textbf{Exact-range factorial} ($50.9$M; $2$ bases $\times$ $2$
training lengths $\times$ $3$ head dimensions $={}12$ configurations, $3$ seeds).
All arms share sampled extrema and log-span with Geo. Metric: weighted OOD NLL
over $2\times/4\times/8\times$, paired within configuration and seed. All arms
were pre-specified; no contrast was registered as primary. CIs are $10{,}000$
configuration-level bootstraps; $p$ is a two-sided exact sign-flip over all
$2^{12}$ assignments, unadjusted across contrasts, so the claim rests on the
cross-configuration direction
(App.~\ref{sec:identification-details}).}
\label{tab:m4}
\begin{tabular}{@{}lcccc@{}}
\toprule
Contrast & $\Delta$NLL & 95\% CI & $\Delta{<}0$ configs & $p$ \\
\midrule
Cosh $0.75\times$ $-$ Geo & $-0.00912$ & $[-0.018,-0.001]$ & $8/12$ & $0.082$ \\
Cosh rule $-$ Geo & $-0.00988$ & $[-0.021,0.002]$ & $7/12$ & $0.125$ \\
Cosh $1.25\times$ $-$ Geo & $\mathbf{-0.01210}$ & $[-0.021,-0.003]$ & $\mathbf{10/12}$ & $0.027$ \\
Matched exponential $-$ Geo & $-0.01062$ & $[-0.021,-0.001]$ & $9/12$ & $0.071$ \\
\midrule
\multicolumn{5}{@{}l@{}}{\textit{The matched exponential reproduces the non-uniform direction:}}\\
Cosh (rule) $-$ exponential & $+0.00074$ & $[-0.006,0.008]$ & $7/12$ & $0.836$ \\
\bottomrule
\end{tabular}
\end{table}


\paragraph{Absolute values.} Mean weighted OOD NLL (equivalent PPL): Geo
$6.211973$ ($498.68$); cosh $0.75\times$ $6.202858$ ($494.16$); cosh
$1.00\times$ $6.202094$ ($493.78$); cosh $1.25\times$ $6.199873$ ($492.69$);
matched exponential $6.201354$ ($493.42$).

\paragraph{Boundary-$\tau$ arms.} At formula $\tau{=}1$ the $1.5\times$ arm is
better; at formula $\tau{=}8$ the $0.75\times$ arm is. Optima move inward from
both ends, which bounds the formula's implied range but does not yield a
replacement closed-form rule.

\paragraph{Scope.} This is a short-budget $50.9$M mechanistic NLL control
trained on a deterministically repeated WikiText-2 raw stream, a different
corpus from the $151.9$M control and from the mature-scale runs. The factorial
establishes cross-configuration \emph{direction}, while the $151.9$M control
provides the effect size within its configuration. Scale is evaluated in
\S\ref{sec:exp-mature}.

\section{The YaRN-style range scaler}
\label{sec:ramp-scaler}

Every arm marked $+\rs{}$ uses the repository's fixed-index YaRN-style range
operator, inspired by the NTK-by-parts/YaRN family \citep{peng2024yarn}:
high-frequency channels are left nearly unchanged and lower-frequency channels
receive progressively more range scaling. It is held
\emph{identical} across the arms being compared --- same implementation, same
fixed cut indices, same scale $s$ --- so that the only difference between
$\mathrm{Geo}+\rs{}$ and $\evq{}+\rs{}$ is the training-time table each
substrate presents to the same operator. It therefore measures
substrate-dependent leverage at the fixed tested scale.

\paragraph{Implemented components.} For $K$ frequency channels, the scaler
uses fixed channel-index cut points
$k_{\mathrm{lo}}=\lfloor0.20K\rfloor$ and
$k_{\mathrm{hi}}=\lfloor0.90K\rfloor$. Between them it applies the cubic
smoothstep $r(x)=x^2(3-2x)$ and replaces each inverse frequency $\omega_k$ by
\[
\omega_k\big/\left(s^{r_k}t^{r_k/2}\right),
\qquad t=1+0.07\log_2 s.
\]
Its boundaries are fixed by channel index; the factor $t$ is folded into
frequency scaling. These implementation choices define \rs{} and distinguish
it from the cited YaRN reference method.

\section{Learned-frequency comparator}
\label{sec:pe-dominant}

\begin{table}[h]
\caption{Learned-frequency comparator ($125$M, $\Ltr{=}128$, FineWeb-Edu,
$128\to8$K). Geo, learned-inv-freq, and \evq{} are seed $42$; learnable-$\tau$
is $3$-seed (42/137/256). Allocation-shape attribution uses the fixed,
zero-parameter analytic schedules of \S\ref{sec:exp-identify}; this separately
reported arm is a layer-shared learnable
inverse-frequency baseline --- one learned scalar per band, shared across
layers and heads --- and therefore belongs to the learned-\emph{table} family
\citep{karypis2026lerope}; DAPE belongs to the learned positional-\emph{operator}
family \citep{zheng2024dape}.}
\label{tab:pe-dominant}
\centering
\small
\setlength{\tabcolsep}{5pt}
\renewcommand{\arraystretch}{0.92}
\begin{tabular}{@{}lcccc@{}}
\toprule
Method ($128\!\rightarrow\!8$K) & Extra params & PPL@$128$ & PPL@$8$K & $\Delta$ vs Geo \\
\midrule
Geo & 0 & 184.9 & 513.7 & --- \\
Learnable $\tau$ & 1 & $\mathbf{181.2{\scriptstyle\,\pm 1.3}}$ & $437.9{\scriptstyle\,\pm 12.2}$ & $-14.8\%$ \\
Learned inv-freq ($100\times$ positional LR) & 32 & 183.6 & 455.3 & $-11.4\%$ \\
\quad same, $10\times$ positional LR & 32 & --- & 477.7 & $-7.0\%$ \\
\evq{} & 0 & 182.0 & \textbf{333.7} & $\mathbf{-35.0\%}$ \\
\bottomrule
\end{tabular}
\end{table}


The $32$-parameter arm in Table~\ref{tab:pe-dominant} is a layer-shared
learnable inverse-frequency baseline: one learned scalar per frequency band,
shared across layers and heads. It received a dedicated positional-parameter
learning-rate sweep at $10\times$ and $100\times$ the base rate, and the
stronger $100\times$ setting is the $455.3$ row. It answers the learned-table
tuning question. Allocation-shape attribution uses the fixed analytic schedules
under an identical zero-parameter operator in \S\ref{sec:exp-identify}.

On method identity: DAPE \citep{zheng2024dape} is a learned positional
\emph{operator}, whereas this comparator learns inverse frequencies directly
and belongs to the learned-\emph{table} family later studied at scale by
\citet{karypis2026lerope}. The distinction matters for what the row can be read
as evidence for. Read as a learned-table comparison, its direction is
consistent with the independent finding that learned tables degrade more
sharply than standard RoPE under naive extrapolation.
